% Section 1: Introduction
\label{sec:introduction}

Artificial intelligence is rapidly becoming the infrastructure of financial markets. Machine learning methods determine which securities are cheap and which are expensive. Large language models produce the financial analyses that inform investment committees. Algorithmic systems calibrate the models that market makers use to price liquidity. This transformation is well documented: diverse machine learning methods converge on the same dominant predictors of equity returns, and large language models at fixed temperature generate near-identical financial statement analyses. The convergence is not an aberration. It is a structural feature of training on common data and deploying pre-trained models at scale.

The central question is: what happens to financial market stability when the private signals of many agents become highly correlated because they are all processed by the same AI system? The question is not about whether AI makes individual agents better informed. It is about what happens when it makes them all informed in the same way.

We formalise AI signal homogeneity through a single primitive: the cross-sectional correlation $\rho \in [0,1]$ of AI-generated signals. Each agent's signal error decomposes into a common AI model error shared by all AI-equipped agents and an idiosyncratic component, with $\rho$ governing the fraction of variance attributable to the common component. At $\rho = 0$, signals are independent and the economy sits at the Morris-Shin heterogeneous-signal benchmark. At $\rho = 1$, every AI-equipped agent observes the same signal. The entire paper is organised as comparative statics in $\rho$.

We show that rising $\rho$ activates fragility through three channels, each connected to an established literature but none sufficient on its own.

Channel 1 (coordination failure) shows that correlated AI signals function as an endogenous public signal in a \citet{goldstein2005} bank-run game, restoring the equilibrium multiplicity of \citet{hellwig2002} and making self-fulfilling crises possible. Channel 2 (information acquisition collapse) shows that cheap but correlated AI signals crowd out costly private research in a \citet{grossman1980} framework, degrading the novel information that prices reveal even as forecast accuracy improves. Channel 3 (correlated liquidity withdrawal) shows that shared AI calibrations reduce the effective number of independent market makers, creating a threshold above which all withdraw simultaneously. Taken individually, each channel has its own fragility threshold. Taken together, they have a lower one.

The paper's central result is that the three channels interact and amplify. We define a composite operator mapping the state vector of effective correlation, crisis threshold, and effective market-maker count into itself through three equilibrium mappings that formalise the cross-channel feedback. By Brouwer's fixed-point theorem, this operator admits at least one fixed point on a compact convex domain.

This fixed point is where the danger lies. Let $\rho_1^*, \rho_2^*, \rho_3^*$ denote the individual channel fragility thresholds and $\rho^*$ the bifurcation threshold of the integrated system. We show that
\begin{equation}\label{eq:main-result-intro}
  \rho^* < \min(\rho_1^*, \rho_2^*, \rho_3^*).
\end{equation}
The integrated system becomes fragile at a strictly lower AI signal correlation than any individual channel assessment predicts. There exists a non-empty interval $(\rho^*, \min_i \rho_i^*)$ in which every single-channel stress test returns a favourable verdict while the system has already crossed its compound fragility threshold. We call this the \emph{safety illusion}. It is the paper's central message: any regulator examining one channel at a time will miss the compound threshold entirely.

The model delivers three main results with direct implications for financial regulation. The crisis threshold in a bank-run game is non-monotone in $\rho$: moderate AI correlation improves individual inference and stabilises the system, but high correlation coordinates withdrawal decisions and makes self-fulfilling crises possible, implying that AI adoption can shift a banking system from stable to fragile without any change in fundamentals. Standard forecast-accuracy metrics paint a misleadingly favourable picture of AI-assisted markets, because AI improves price prediction while degrading the novel information that prices reveal to firm managers, so regulators relying on forecast accuracy will miss a deterioration in the real-economy information channel. The effective number of independent liquidity providers collapses rapidly with AI model similarity, and a diversity mandate bounding institutional AI model similarity below the compound threshold $\rho^*$ corrects the resulting prisoner's dilemma in AI adoption.

The closest predecessors are \citet{danielsson2025} and \citet{yang2024}. \citet{danielsson2025} derive a crisis threshold as a function of $\mu$, the fraction of AI agents, through an execution-speed channel. Their primitive captures the extensive margin of AI penetration; ours captures the intensive margin of AI output homogeneity. Moreover, \citet{danielsson2025} formally model only the coordination/run channel and do not integrate the \citet{hellwig2002} multiplicity result into their formal model; they do not model information acquisition collapse, market-maker liquidity withdrawal, or the cross-channel amplification loop. \citet{yang2024} demonstrates AI coordination computationally via Q-learning agents but has no signal correlation parameter and no analytical characterisation. \citet{hansen2025} find that independently queried AI agents reduce herding, consistent with the low-$\rho$ regime; they do not test the high-$\rho$ regime where our fragility mechanism activates.

In the information acquisition literature, \citet{dugast2018} and \citet{farboodi2020} model temporal tradeoffs in data technology but do not address the cross-sectional dimension. In market microstructure, \citet{danielsson2012} formalise procyclical selling from common VaR constraints, which we nest as the $\rho = 1$ limiting case. \citet{danielsson2022} provide the most comprehensive qualitative argument for multi-channel AI systemic risk; \citet{dou2025a, dou2025b} model AI fragility through reinforcement learning. Our mechanism is complementary, arising from signal homogeneity rather than strategic learning. Section~\ref{sec:literature} provides a full discussion.

Section~\ref{sec:model} introduces the signal structure, agent types, timing, and equilibrium concepts. Sections~\ref{sec:channel1}--\ref{sec:channel3} analyse Channels 1, 2, and 3 respectively. Each section demonstrates that its channel creates fragility in isolation; each also reveals why that channel alone provides an incomplete picture. Section~\ref{sec:amplification} formalises the amplification loop and derives the bifurcation result. Section~\ref{sec:extensions} endogenises $\rho$ and characterises the diversity mandate. Section~\ref{sec:empirics} provides motivating evidence using the ChatGPT release as an AI adoption shock. Section~\ref{sec:conclusion} discusses policy implications and future work.
% --- Clarity pass complete: 2026-03-10 ---
% --- Flow pass complete: 2026-03-11 ---
% --- Technical audit complete: 2026-03-11 ---
